\phantomsection\label{fig:vol01:spring-autumn:qi-huan-alliance-network}
\begin{center}
\begin{tikzpicture}[x=.88cm,y=.88cm,font=\small]
  \fill[paper] (0,0) rectangle (11.4,6.0);

  \fill[dynasty!20] (8.15,3.45) -- (10.25,3.4) -- (10.55,5.15) -- (8.4,5.3) -- cycle;
  \node[align=center] at (9.35,4.35) {齐\\会盟的组织者};
  \fill[source!16] (3.95,4.7) -- (5.6,4.65) -- (5.75,5.6) -- (4.1,5.7) -- cycle;
  \node[align=center] at (4.85,5.15) {周王室\\名分与礼仪};
  \fill[timeline!18] (6.25,2.65) -- (7.65,2.6) -- (7.85,3.7) -- (6.45,3.8) -- cycle;
  \node[align=center] at (7.05,3.2) {葵丘\\前651年};
  \fill[source!16] (5.95,5.25) -- (8.0,5.2) -- (8.15,5.9) -- (6.15,5.95) -- cycle;
  \node[align=center] at (7.05,5.58) {邢、卫等\\北方援助关系};
  \fill[timeline!18] (2.15,.75) -- (3.85,.7) -- (4.05,1.85) -- (2.35,1.95) -- cycle;
  \node[align=center] at (3.05,1.3) {召陵\\前656年};
  \fill[source!15] (.65,1.8) -- (1.95,1.75) -- (2.1,3.0) -- (.85,3.1) -- cycle;
  \node[align=center] at (1.4,2.4) {楚\\江汉方向};
  \fill[timeline!14] (4.3,2.05) -- (5.4,2.0) -- (5.55,2.95) -- (4.45,3.05) -- cycle;
  \node[align=center] at (4.9,2.5) {宋};
  \fill[timeline!14] (5.0,3.6) -- (5.9,3.55) -- (6.05,4.35) -- (5.1,4.45) -- cycle;
  \node at (5.5,4.0) {鲁};

  \draw[-{Stealth[length=2mm]},ultra thick,color=dynasty]
    (8.35,4.35) -- (7.8,3.55);
  \draw[-{Stealth[length=2mm]},thick,color=dynasty]
    (8.25,4.0) .. controls (6.1,2.0) and (4.75,1.3) .. (4.05,1.35);
  \draw[-{Stealth[length=2mm]},thick,color=timeline]
    (8.65,5.1) -- (7.85,5.5);
  \draw[-{Stealth[length=2mm]},thick,color=timeline]
    (5.1,2.8) -- (6.2,3.05);
  \draw[-{Stealth[length=2mm]},thick,color=timeline]
    (5.7,3.9) -- (6.35,3.5);
  \draw[{Stealth[length=2mm]}-{Stealth[length=2mm]},thick,dashed,color=source]
    (2.15,1.95) .. controls (1.75,2.25) and (1.75,2.5) .. (2.05,2.55);

  \node[text=dynasty,align=center] at (5.2,1.12) {齐率盟军南向\\（约略方向）};
  \node[text=source,align=center] at (1.95,3.35) {对峙与交涉\\非臣属线};
  \node[text=timeline,align=center] at (8.65,2.5) {会盟以诸侯\\即时承认为条件};
  \node[text=source,align=center,font=\footnotesize] at (9.25,.85)
    {前659、656、651年是连续的外交节点，\\不是一张同时存在的固定联盟地图};
\end{tikzpicture}

\smallskip
\small\textit{图\quad 齐桓公时期的会盟网络示意：从北方援助、召陵对楚交涉到葵丘会盟，显示齐如何把军事援助、周室名分与诸侯承认连在一起。参考《春秋·僖公元年、四年、九年》《左传·僖公元年、四年、九年》《国语·齐语》及《史记·齐太公世家》；盟员、会期和动机在传世叙事中仍有可讨论处。国家范围、会盟地点与往来路线均为约略示意，不能视作精确疆界、行军图或固定从属关系。}
\end{center}
